\section{Описание}
Требуется реализовать алгоритм точного поиска подстроки в тексте с использованием Z-блоков (Z-функции). 

Как описано в \cite{Gasfield}, Z-функция от строки $S$ представляет собой массив $Z$, где $Z[i]$ равно длине наибольшего общего префикса строки $S$ и её суффикса, начинающегося с позиции $i$. Главное преимущество алгоритма заключается в том, что он строит массив $Z$ за линейное время $O(N)$. 

Для решения задачи поиска образца $P$ в тексте $T$, строки объединяются через специальный символ-разделитель (отсутствующий в алфавите) в новую строку $S = P + \# + T$. Если в вычисленном массиве $Z$ для позиции $i$ выполняется условие $Z[i] = |P|$, это означает успешное нахождение образца.

\textbf{Особенности реализации:}
В данной работе алгоритм адаптирован для поиска не по отдельным символам, а по целым словам:
\begin{enumerate}
\item Вместо строки символов используется вектор строк \texttt{std::vector<std::string>}. Сравнение происходит по целым словам за $O(L)$, где $L \le 16$ — максимальная длина слова.
\item Для обеспечения регистронезависимости символы приводятся к нижнему регистру с помощью быстрой ASCII-арифметики в момент посимвольного чтения \texttt{std::cin.get()}, что позволяет избежать выделения дополнительной памяти.
\item Для правильного вывода номеров строк и слов в тексте параллельно с массивом слов заполняется массив метаданных \texttt{std::vector<WordPos>}.
\end{enumerate}

\section{Исходный код}
Основная логика алгоритма заключена в функции \texttt{computeZ}, работающей с вектором строк. 

\begin{lstlisting}[language=C++]
std::vector<int> computeZ(const std::vector<std::string>& words) {
    int n = (int)words.size();
    std::vector<int> z(n, 0);
    int l = 0, r = 0;
    for (int i = 1; i < n; i++) {
        if (i <= r) {
            z[i] = std::min(r - i + 1, z[i - l]);
        }
        while (i + z[i] < n && words[z[i]] == words[i + z[i]]) {
            z[i]++;
        }
        if (i + z[i] - 1 > r) {
            l = i;
            r = i + z[i] - 1;
        }
    }
    return z;
}
\end{lstlisting}

Полный листинг программы (включая парсинг текста) приведен в приложении:

\lstinputlisting[language=C++, caption={Поиск подстроки с помощью Z-функции}, label={lst:maincpp}]{../main.cpp}
\pagebreak

\section{Консоль}
Пример сборки программы и тестирования на данных из условия.

\begin{alltt}
$ g++ -std=c++20 -Wall -Wextra -O2 main.cpp -o z_search
$ cat test.in
cat dog cat dog bird
CAT dog CaT Dog Cat DOG bird CAT
dog cat dog bird
$ ./z_search < test.in > result.out
$ cat result.out
1, 3
1, 8
\end{alltt}
\pagebreak